%% kapa: exam=1 round=29 year=2018 period=1 form=A subject=economics subject_ko=경제학원론 seq=40 number=80
%% source: Q-Net 감정평가사 기출문제 (cst00309) artlSeq=5202808 file=extracted/1008/5202808/2018년 제29회 감정평가사 1차시험 기출문제/2018년 제29회 감정평가사 1차시험 기출문제/1교시 A형.pdf
%% conversion: kordoc 4.24.51 -> tools/gen1.py (auto); review=proofread-sonnet; flags=none
\begin{question}{80}
경기안정화 정책에 관한 설명으로 옳은 것은?
\begin{choices}
  \item 재정지출 증가로 이자율이 상승하지 않으면 구축효과는 크게 나타난다.
  \item 투자가 이자율에 비탄력적일수록 구축효과는 크게 나타난다.
  \item 한계소비성향이 클수록 정부지출의 국민소득 증대효과는 작게 나타난다.
  \item 소득이 증가할 때 수입재 수요가 크게 증가할수록 정부지출의 국민소득 증대효과는 크게 나타난다.
  \item 소득세가 비례세보다는 정액세일 경우에 정부지출의 국민소득 증대효과는 크게 나타난다.
\end{choices}
\end{question}
\answer{5}
